% 01 Introduction
\section{引言}
\label{sec:introduction}

% 序列数据广泛存在于工程系统与科学观测中，是刻画复杂动态过程与观测机制的核心载体。时间序列记录了系统随时间变化的观测值，因而在制造、物流、能源与医疗保健等多个领域至关重要：例如预测电力消耗、交通流量、产品销量与库存水平，或对工业过程进行早期预警与资源调度。在科学观测中，序列数据同样扮演关键角色：诸如由光谱/测光等观测序列反演天体的物理与化学属性，可为恒星演化与银河化学演化研究提供基础支撑，其中 CEMP（Carbon-Enhanced Metal-Poor）恒星的参数回归是具有代表性的应用场景。从学习形式上看，这类问题可统一表述为\emph{序列到多变量输出}学习：既包括多变量时间序列预测（MTSF）的多步预测（sequence-to-sequence），也包括参数反演的多变量回归（sequence-to-vector）。尽管输出形态不同，两者都需要在噪声与观测不完备条件下刻画跨维耦合与长程依赖，并在分布漂移与统计量变化下保持稳定泛化。

序列数据广泛存在于工程系统与科学观测中，是刻画复杂动态过程与观测机制的核心载体。时间序列记录了系统随时间变化的观测值，因而在制造、物流、能源与医疗保健等多个领域至关重要：例如预测电力消耗、交通流量、产品销量与库存水平，或对工业过程进行早期预警与资源调度。在科学观测中，序列数据同样扮演关键角色：诸如由光谱/测光等观测序列反演天体的物理与化学属性，可为恒星演化与银河化学演化研究提供基础支撑，其中 CEMP（Carbon-Enhanced Metal-Poor）恒星的参数回归是具有代表性的应用场景。尽管上述任务在目标形式上分别对应未来序列生成与物理参数估计，但其本质都属于基于序列观测的多变量建模问题，即模型不仅需要从复杂时序结构中捕获跨变量耦合关系与长程依赖，还需在噪声干扰、观测不完备以及数据分布变化等条件下保持稳定的泛化能力。

基于从传统统计模型到深度模型（如 RNN/TCN），特别是近年来基于 MLP 和 Transformer 架构的研究进展，加上输入表示层面（如 Patch 分块与协变量引入）的持续优化，已经迅速提高了时间序列预测的精度。然而，现有的深度网络在处理复杂多变量任务时，往往依赖隐式的频率状态利用和单一的结构先验，使其在非平稳条件下的动态适应能力严重受限。这种缺乏协同动态调制与稳定化机制的缺陷在实际应用中是具有不确定性的。在具有高噪声和异质数据的现实场景中，系统不太可能依赖那些难以应对分布漂移且极易出现训练不稳的预测模型。在对泛化稳定性和可靠性至关重要的领域部署此类鲁棒性不足的模型会带来重大风险。如果能在高精度的同时显式引入与时序信号一致的频率归纳偏置并协同跨变量建模，模型就能够有效应对极端变化。这一主题被称为非平稳时序的鲁棒建模，并在复杂场景的时间序列预测与参数回归领域被积极研究。一些现有方法试图通过降低长序列注意力复杂度（Informer~\cite{informer}）、从分解与频域角度增强周期结构刻画（Autoformer~\cite{autoformer}; FEDformer~\cite{fedformer}）、通过变量维 token 化提升跨变量交互能力（iTransformer~\cite{itransformer}）以及针对分布漂移提供可逆归一化机制（RevIN~\cite{revin}）来应对这些挑战。然而，如何在高精度的同时引入与时序信号一致的频率归纳偏置，并在非平稳条件下保持训练与泛化稳定性，仍是模型走向实际应用的关键问题。

为解决上述问题，本文提出\textbf{渐进式频率条件 FiLM}（Progressive Frequency-conditioned FiLM），一种以 Transformer 为骨干、通过逐层渐进衰减的 FiLM 将多分辨率频谱信息注入 Transformer Encoder 的轻量高效框架。其核心创新包含三个要素：

\textbf{(1) 多分辨率频谱提取（Multi-Resolution Spectrum）}：在三个时间窗口（全序列 \(L\)、半序列 \(L/2\)、四分之一序列 \(L/4\)）上分别通过 FFT 提取功率谱，每窗口独立计算多频带归一化能量、高低频比值与频谱熵，形成尺度鲁棒且信息丰富的频域表征。相比固定低通滤波，FFT 频谱能够更完整地捕获周期结构与谱平坦度等可预测性信号。

\textbf{(2) 逐层渐进式 FiLM（Layer-wise Progressive FiLM）}：压缩后的频谱 latent 通过每层独立的小 MP 生成调制参数 \((\alpha, \beta)\)，注入到每一个 Transformer Encoder Layer 之前；调制强度按 \(\gamma_{\text{base}} \cdot \delta^{\ell}\) 随层深指数衰减（\(\delta \in (0,1)\) 为衰减因子），实现由粗到细的层级化频率条件调控——浅层接受较强的频域引导以建立全局谱结构感知，深层逐渐减弱调制以保护语义表征的精细辨别能力。

\textbf{(3) 可学习调制强度（Learnable Modulation Scale）}：将全局调制基强度 \(\gamma_{\text{base}}\) 参数化为 \(\mathrm{sigmoid}(\theta_{\gamma})\)（其中 \(\theta_{\gamma}\) 为可学习参数，初始对应 \(\gamma_{\text{base}} \approx 0.02\)），替代手工设定的固定系数。不同数据集在训练中自适应决定频域信息的注入程度，同时兼顾可退化性：初始状态下调制接近关闭，模型等价于纯 iTransformer。

上述设计与 iTransformer~\cite{itransformer} 的变量维 token 化策略及 RevIN~\cite{revin} 可逆归一化无缝协同，并通过 FiLM 末层零初始化、\(\tanh\) 有界约束与 \(\log(1+x)\) 压缩等稳定化设计保证训练稳健。频率条件信号具有明确物理/信号含义（频带能量分布、高低频能量对比、频谱平坦度），为理解模型在不同序列形态下的响应差异提供了可分析性。框架既可用于 MTSF 多步预测，也可通过更换任务头用于 CEMP 等场景下的多变量参数回归。

\subsection{本文贡献}
\label{subsec:contributions}
本文的主要贡献概括如下（可根据最终实验进一步精炼/增补）：
\begin{itemize}
  \item 提出\textbf{多分辨率频谱特征提取}方法，以 FFT 功率谱的多频带归一化能量、高低频比值与频谱熵联合刻画输入谱状态，相比传统低通滤波方案具有更强的尺度鲁棒性与信息完整性。
  \item 设计\textbf{逐层渐进式 FiLM 调制}机制：每层独立生成调制参数并以指数衰减强度逐层注入，在不破坏 iTransformer 主干可训练性的前提下实现层级化频域引导。
  \item 提出\textbf{可学习调制强度}参数化方案（Sigmoid 门控 + 保守初始化），使模型自适应决定频域信息的注入程度，在无频域信号的数据上可自动退化为纯 iTransformer。
  \item 给出\textbf{面向非平稳与训练稳定的整体方案}：RevIN 处理分布漂移，配合频谱压缩、\(\log(1+x)\) 压缩/截断、\(\tanh\) 有界调制与零初始化，增强鲁棒性与可训练性。
  \item 在\textbf{多变量预测与 CEMP 参数回归}两类任务上的实验与消融分析（待补充）验证方法有效性，并量化各模块对性能的贡献。
\end{itemize}

\subsection{论文结构}
本文余下部分组织如下：第\ref{sec:related_work}节介绍背景与相关工作；
第\ref{sec:method}节给出方法细节；
第\ref{sec:experiments}节描述实验设置并报告结果；
第\ref{sec:conclusion}节总结全文并讨论未来工作。
